% Part: turing-machines
% Chapter: machines-computations
% Section: halting-states

\documentclass[../../../include/open-logic-section]{subfiles}

\begin{document}

\olfileid{tur}{mac}{hal}
\olsection{درېدني حالتونه}

\begin{explain}
که څه هم موږ خپل ماشينونه داسې تعريف کړي چې يوازې د عملي کېدونکې
لارښوونې د نشتوالي پر وخت درېږي، د ټيورينګ ماشينونو په عامو
څرګندونو کښې يو ځانګړی \emph{درېدنی حالت}~$h$ وي، داسې چې
$h \in Q$۔

د درېدني حالت تر شا نظر ساده دے: کله چې ماشين خپل کار بشپړ کړي (د
ننوت منلو ته چمتو وي، يا د وت ليکل ئې بشپړ کړي وي)، داسې حالت~$h$
ته ننوځي چې هلته درېږي۔ ځينې ماشينونه دوه درېدني حالتونه لري: يو
ننوت مني او بل ننوت ردوي۔
\end{explain}

\begin{ex}\emph{درېدني حالتونه}۔
د دې مفهوم د روښانولو دپاره راځئ د جفت ماشين له يو بدلون څخه پيل
وکړو۔ د دې پر ځاے چې ماشين د جفت ننوت په صورت کښې په حالت~$q_0$
کښې ودرېږي، يوه لارښوونه ور زياتولے شو چې ماشين درېدني حالت ته
واستوي۔
\[
\begin{tikzpicture}[->,>=stealth',shorten >=1pt,auto,node distance=2.8cm,
                    semithick]
  \tikzstyle{every state}=[fill=none,draw=black,text=black]

  \node[initial, state]         (A)                     {$q_0$};
  \node[state]         (B) [right of=A] {$q_1$};
  \node[state]         (C) [below of=A] {$h$};

  \path (A) edge [bend left] node {\TMtrans{\TMstroke}{\TMstroke}{R}} (B)
  	    edge node {\TMtrans{\TMblank}{\TMblank}{N}} (C)
        (B) edge [loop above] node {\TMtrans{\TMblank}{\TMblank}{R}} (B)
            edge [bend left] node {\TMtrans{\TMstroke}{\TMstroke}{R}} (A);
\end{tikzpicture}
\]

راځئ بېلګه نوره هم پراخه کړو۔ کله چې ماشين ومومي چې ننوت طاق دے،
هېڅکله نۀ درېږي۔ موږ د کړۍ وهونکې لارښوونې پر ځاے داسې لارښوونه
ايښودلے شو چې رد حالت~$r$ ته ځي، او په دې توګه ماشين ته يو
\emph{رد} حالت ور زياتولے شو۔
\[
\begin{tikzpicture}[->,>=stealth',shorten >=1pt,auto,node distance=2.8cm,
                    semithick]
  \tikzstyle{every state}=[fill=none,draw=black,text=black]

  \node[initial,state]         (A)                     {$q_0$};
  \node[state]         (B) [right of=A] {$q_1$};
  \node[state]         (C) [below of=A] {$h$};
  \node[state]         (D) [below of=B] {$r$};

  \path (A) edge [bend left] node {\TMtrans{\TMstroke}{\TMstroke}{R}} (B)
  	    edge node {\TMtrans{\TMblank}{\TMblank}{N}} (C)
        (B) edge node {\TMtrans{\TMblank}{\TMblank}{N}} (D)
            edge [bend left] node {\TMtrans{\TMstroke}{\TMstroke}{R}} (A);
\end{tikzpicture}
\]
\end{ex}

\begin{explain}
د ځانګړي درېدني حالت زياتول د دې په څېر حالتونو کښې ګټور وي، ځکه په
څرګنده ښيي چې ماشين ټاکلي ننوتونه کله مني يا ردوي۔ خو پام کول مهم دي
چې د ځانګړي درېدني حالت په زياتولو هېڅ محاسبوي ځواک نۀ زياتېږي۔ همداراز
د درېدو لږ صوري مفهوم خپلې ګټې لري۔ د درېدو هغه تعريف چې تر اوسه مو
په دې څپرکي کښې کارولے، د \emph{درېدو مسئلې} ثبوت شهودي او څرګندول ئې
اسانه کوي۔ له همدې امله موږ خپل اصلي تعريف ته دوام ورکوو۔
\end{explain}

\end{document}
